% PreTextuais/Resumo.tex
\begin{resumo}
A gestão de frotas no setor público historicamente opera sob rotinas reativas emergenciais devido à fragmentação de dados entre cadastros patrimoniais, ordens de serviço (OSs) e abastecimentos. Conduzida sob os preceitos do método lógico para redação científica (Volpato, 2015, 2017), esta dissertação propõe e avalia uma arquitetura computacional de analítica prescritiva orientada por Inteligência Artificial Explicável (XAI), Teoria dos Grafos e Otimização Inteira Mista (MILP), estruturando-se pela prevalência do conteúdo sobre a forma e pela distinção categórica entre significância estatística e relevância prática operacional. A base empírica abrange 77.867 ordens de serviço transacionadas entre 2020 e 2026 no Poder Executivo de Santa Catarina, totalizando R\$ 138,76 milhões em dispêndio financeiro aprovado e liquidado (após a retenção de R\$ 60,05 milhões em glosas técnicas preventivas pela fiscalização contratual), vinculadas a 34.280 veículos e 32,84 milhões de litros de combustível. As hipóteses foram avaliadas pelo texto do anteprojeto e pelo limiar numérico pré-registrado. No Cenário S0 (2.000 veículos, 36 meses, instância semissintética), a política MILP resolvida pelo HiGHS custou 13,66\% (R\$ 6,99 milhões) menos que a prática reconstruída, com todas as manutenções obrigatórias cumpridas (H3 sustentada), embora uma heurística gulosa alcance praticamente o mesmo custo. A integração do patrimônio às ordens de serviço (H1) melhorou a ROC-AUC em +0,0095 (IC 95\% 0,0045 a 0,0146), mas reduziu o erro médio absoluto em apenas 0,62\%, abaixo da meta de 15\% (parcialmente sustentada). Os ensembles de árvores superaram o SVC-RBF na classificação, mas erraram 8,65\% a mais que o SVR-RBF na regressão (H2 não sustentada). O SHAP mostrou-se exato e estável, mas o LIME teve baixa fidelidade local (R\textsuperscript{2} = 0,243) e a concordância entre os métodos ficou em Jaccard 45,48\% (H4 não sustentada). Identificaram-se 242 veículos com gasto acumulado superior a 40\% do valor contábil, como triagem para decisões de alienação. Conclui-se que a analítica prescritiva pode reduzir o custo esperado da manutenção sob as mesmas restrições de orçamento, e que a explicabilidade para auditoria deve se apoiar em métodos fiéis ao modelo.

\vspace{\onelineskip}
\noindent
\textbf{Palavras-chave:} analítica prescritiva; gestão de frotas públicas; teoria dos grafos; inteligência artificial explicável; programação linear inteira mista.
\end{resumo}
